\chapter*{第一篇勘误与评论}
\addcontentsline{toc}{chapter}{第一篇勘误与评论}
\phantomsection\label{errata-ct}
\noindent\textcolor{Muted}{以下译自原文件附载的 2008 年勘误两页，对应第一篇原书第 1--49 页。所引页码及行号均指原书；负行号从页底向上数。}

\begin{description}
\item[第 14 页，页底]
Chad Schoen（\emph{J. Algebraic Geom.}\ \textbf{11}（2002），41--100）已经证明，这个映射一般并非单射。
\hfill\hyperref[orig:19]{返回正文}

\item[第 18 页，倒数第 5 行]
应为“réduction”（约化）。\hfill\hyperref[orig:23]{返回正文}

\item[第 24 页，第 19 行]
应为“D'après”（根据）。\hfill\hyperref[orig:29]{返回正文}

\item[第 32 页，第 2 行]
应为“apparue”（出现）。\hfill\hyperref[orig:37]{返回正文}

\item[第 38 页，第 24 行]
应引用 [CR3]。\hfill\hyperref[orig:43]{返回正文}

\item[第 40 页，定理 8.4 及备注 8.4.1]
定理 8.4 中关于 Albanese 簇具有好约化的假设不能删去；R. Parimala 与 V. Suresh 的一个例子说明了这一点，见 \emph{Inventiones math.}\ \textbf{122}（1995），83--117。
\hfill\hyperref[orig:45]{返回正文}

\item[第 41 页，第 6 行]
应读作
\[
 H^2\!\left(K,H^1_{\et}\!\left(\overline X,\mathbb Q/\mathbb Z(2)\right)^0\right).
\]
\hfill\hyperref[orig:46]{返回正文}

\item[第 42 页，倒数第 5 行至倒数第 1 行]
在定理 8.6 中补充假设：$X/R$ 的相对维数为 $2$。
\hfill\hyperref[orig:47]{返回正文}

\item[第 45 页，第 10 行]
删去一个“le”。\hfill\hyperref[orig:50]{返回正文}

\item[第 45 页，倒数第 2 行]
补充：利用从 $L$ 到 $k$ 的迹，可以看出映射
\[
 H^1(X,\mathcal K_2)\otimes\mathbb Q
 \longrightarrow
 \bigoplus_{p\in\Spec A^{(1)}}\Pic(\mathbb X_p)\otimes\mathbb Q
\]
是满射。\hfill\hyperref[orig:50]{返回正文}

\item[第 46 页，第 2 行]
将
\[
 CH_2(\mathbb X_B)_{\tors}\longrightarrow CH^2(X_L)_{\tors}
\]
替换为
\[
 CH_2(\mathbb X)_{\tors}\longrightarrow CH^2(X)_{\tors}.
\]
勘误页的左端仍印为下标 $2$；正文依据余维二的上下文采用上标 $2$，并另加译注。\hfill\hyperref[orig:51]{返回正文}

\item[第 46 页，倒数第 14 行]
应为“exacte”（正合）。\hfill\hyperref[orig:51]{返回正文}

\item[第 46 页，倒数第 12 行]
将 $H^i(\mathbb X,\mathbb Z_l(j))$ 换成该群模去其挠子群所得的商群。
\hfill\hyperref[orig:51]{返回正文}

\item[第 46 页，倒数第 6 行]
删去“les”。\hfill\hyperref[orig:51]{返回正文}

\item[第 48 页，参考文献 [Sb2]]
Per Salberger 的文章已经发表：\emph{Torsion cycles of codimension 2 and $l$-adic realizations of motivic cohomology}，载于 \emph{Séminaire de théorie des nombres, Paris 1991--1992}（Sinnou David 编），Progress in Mathematics \textbf{116}，Birkhäuser（1993），第 247--277 页。
\hfill\hyperref[orig:53]{返回正文}
\end{description}

\section*{补充参考文献}
\begin{enumerate}
\item J.-L. Colliot-Thélène，\emph{L'arithmétique des zéro-cycles}（1993 年 9 月波尔多算术研讨会报告），\emph{Journal de théorie des nombres de Bordeaux}\ \textbf{7}（1995），51--73。题名译为“零闭链的算术”。
\item J.-L. Colliot-Thélène，\emph{Zéro-cycles sur les surfaces sur un corps $p$-adique : quelques résultats obtenus au moyen de la K-théorie algébrique}，为 Morelia（2003 年 6--7 月）和 Sestri Levante（2004 年 6--7 月）会议准备的讲义，可在原勘误所列网页 \url{http://www.math.u-psud.fr/~colliot/liste-cours-exposes.html} 获取。题名译为“$p$ 进域上曲面的零闭链：借助代数 $K$ 理论得到的若干结果”。
\item Masanori Asakura 与 Shuji Saito，\emph{Surfaces over a $p$-adic field with infinite torsion in the Chow group of $0$-cycles}，\emph{Algebra and Number Theory}\ \textbf{1}（2007），163--181。题名译为“零闭链 Chow 群中含有无限挠子群的 $p$ 进域上曲面”。
\item Shuji Saito 与 Kanetomo Sato（附 Uwe Jannsen 撰写的附录），\emph{A finiteness theorem on zero-cycles over $p$-adic fields}，待刊于 \emph{Annals of Math.}。题名译为“$p$ 进域上零闭链的一个有限性定理”。此处“待刊”沿用勘误的记载。
\item Shuji Saito 与 Kanetomo Sato，\emph{Torsion cycle class maps in codimension two of arithmetic schemes}，预印本 arXiv:math/0612081。题名译为“算术概形余维二的挠闭链类映射”。
\end{enumerate}
